\section{Conclusiones}

La evolución desde métodos estadísticos clásicos hasta redes neuronales de impulsos evidencia un patrón consistente: cada paradigma emergente busca superar las limitaciones anteriores mediante nuevos compromisos entre precisión y eficiencia computacional. Los métodos de aprendizaje profundo demuestran capacidad superior en captura de patrones temporales complejos, pero con costes energéticos considerables; en contraste, las SNNs ofrecen mejoras de eficiencia energética de hasta 3.5 veces respecto a arquitecturas CNN, aunque con ligeras reducciones en precisión.

La arquitectura SNN base bicapa, fundamentada en neuronas LIF y aprendizaje STDP modificado, demuestra viabilidad del enfoque pero expone limitaciones: ausencia de extracción de características jerárquicas y escalabilidad limitada. La optimización bayesiana con Optuna permite cerrar significativamente la brecha de precisión preservando las ventajas energéticas inherentes a las SNNs.

En conclusión, el desarrollo de sistemas eficientes de detección de anomalías requiere arquitecturas híbridas que combinen estratégicamente las fortalezas de múltiples enfoques, particularmente para aplicaciones que demanden simultáneamente alta precisión y eficiencia energética.